\documentclass[10pt,a4paper]{amsart}
\usepackage[margin=19mm]{geometry}
\usepackage{amsmath,amssymb,mathtools}
\usepackage[scr=rsfs,cal=euler]{mathalfa}
\usepackage{xfrac}
\usepackage[unicode,hidelinks,bookmarks=false]{hyperref}
\newcommand{\ie}{i.e.\ }
\newcommand{\cf}{cf.\ }
\newcommand{\resp}{respectively\ }
\newcommand{\Subsection}{\subsection}
\providecommand{\nobreakdash}{\nobreak\hbox{-}}
\setlength{\emergencystretch}{2em}
\multlinegap0pt
\usepackage{mathtools}
\usepackage{amssymb}
\usepackage[shortlabels]{enumitem}
\usepackage{tikz}
\newtheorem{proposition}{Proposition}
\newcommand{\T}[1][]{\mathbb{T}^#1}
\newcommand{\bigPara}[1]{\Bigl(#1\Bigr)}
\newcommand{\la}{\lambda}
\newcommand{\lpn}[2]{||{#1}||_{#2}}
\title{Sharp Eigenfunction Bounds on the Torus for large $p$}
\author{Daniel Pezzi}
\date{Source: arXiv:2603.10927v2 (CC BY 4.0)}
\begin{document}
\maketitle

\noindent\emph{Source and licence.} Sharp Eigenfunction Bounds on the Torus for large $p$. Version arXiv:2603.10927v2 (CC BY 4.0), distributed by arXiv under the Creative Commons Attribution 4.0 International licence, http://creativecommons.org/licenses/by/4.0/, as recorded in the arXiv licence metadata for this exact version and shown on the abstract page https://arxiv.org/abs/2603.10927v2. Full licence notice: \texttt{licenses/arxiv-2603.10927-CC-BY-4.0.txt}.\par\smallskip

\noindent\textit{Editorial scope.} Counted unit: the article's proposition KerrFourierBoundProp (source0.tex lines 357--363). The article's complete original proof of it is delivered with it (source0.tex lines 365--407, one \texttt{proof} environment of 43 lines). No mathematics is written, repaired or re-derived: the statement and the proof are the article's own lines, in the article's own notation, and no retained line is substituted or re-flowed.  Retained uncounted as the source-local context the proof needs: lines 224--226.  Nothing was omitted from any retained span, so the package carries no omission record.  No retained line contains a citation, so the wrapper carries no bibliography and no bibliography entry.  No retained line contains a figure environment, an image-inclusion command or drawing code, so no image is delivered and the package declares no asset dependency.  Editorial formatting only, in addition: an AMS \texttt{amsart} preamble that restates the article's own declarations and macros used by the retained lines, namely the environments \texttt{proposition} on separate counters, together with the standard packages those lines need.  The retained environments print in this document as Proposition 1 in order of appearance, whereas the article prints the same items with the numbering of its own full document.  The complete native source of the article as distributed in the arXiv e-print for this exact version is bundled unchanged as \texttt{sources/196176/source0.tex}.
\begin{equation}\label{2: etaQDef}
    \eta_Q = c_Q \sum_{a/q\in R_{Q}}\eta((t-\frac{a}{q})Q^2).
\end{equation}

\begin{proposition}\label{2: KerrFourierBoundProp}
    For any $d\geq 1$ and $N\leq Q\leq N^2$, we have

    \begin{equation}\label{2: KerrFourierBoundEqn}
        \lpn{\mathcal{F}(K-K^Q)}{\infty}\lesssim \log Q\,Q^{-1}
    \end{equation}
\end{proposition}

\begin{proof}
    We compute the Fourier coefficient of $K-K^Q$ as

    \begin{equation*}
        \mathcal{F}(K-K^Q)(\mathbf{k}) = \int_{\T{d}}\int_{[0,1]}\prod_{j=1}^dG(t,x_j)\bigPara{1-\eta_Q(t)} e(-\la t)e(-\mathbf{k}\cdot x)dtdx.
    \end{equation*}
    After exchanging the order of integration, we see that the integral with respect to $x$ is only non-zero when $\mathbf{k} = (k_1,...,k_d)$ where $k_j$ come from the sum in the definition of $G(t,x_j)$. This then leaves us with only a Fourier transform with respect to time, and so 

    \begin{equation*}
        \mathcal{F}(K-K^Q)(\mathbf{k}) = \widehat{1-\eta_Q}(|\mathbf{k}|^2-\la)\prod_{j=1}^d\gamma(\frac{k_j}{N}).
    \end{equation*}
    Define $l = |\mathbf{k}|^2-\la$. The bump functions do not alter our conclusions as they equal $1$ if $l = 0$. As $1-\eta_Q(t)$ has mean zero, we get

    \begin{equation*}
        \mathcal{F}(K-K^Q)(\mathbf{k}) = 0 \text{ if } |\mathbf{k}|^2=\la.
    \end{equation*}
    If $l\neq 0$, the constant term drops out and we compute the Fourier transform of $\eta_Q$. Recalling \eqref{2: etaQDef}, we have

    \begin{equation}\label{2: FourierTransNumberThrBoundEqn}
        \widehat{1-\eta_Q}(l) \sim c_Q Q^{-2}\hat{\eta}(\frac{l}{10Q^2})\sum_{q\in R_Q}\sum_{a=1}^{q-1}e(la/q).
    \end{equation}
    This follows from basic properties of the Fourier transform. Recall the $c_Q$ factor is $\sim \log Q$. As the sum is over roots of unity, the inner sum is $-1$ if $q$ does not divide $l$ and $q-1$ otherwise. Therefore, we have

    \begin{equation*}
        \widehat{1-\eta_Q}(l) \sim \log Q\, Q^{-2}\hat{\eta}(\frac{l}{10Q^2})\bigPara{-\#D_Q+\sum_{q\in D_Q,\, q|l}q}.
    \end{equation*}
    We have $\#D_Q\sim Q (\log Q)^{-1}$ and $q\sim Q$. Suppose $l$ has less than 10 factors in $D_Q$. Then the triangle equality immediately implies that

    \begin{equation*}
        |\widehat{1-\eta_Q}(l)| \sim \log Q\,Q^{-2}|\hat{\eta}(\frac{l}{10Q^2})|\cdot 10 Q \lesssim \log Q\,Q^{-1}.
    \end{equation*}
    Now suppose $l$ has $S$ factors with $S\geq 10$. Then $l\gtrsim Q^{S}$. However, $\eta$ is Schwartz and so enjoys rapid decay.

    \begin{equation*}
        |\hat{\eta}(z)|\lesssim (1+|z|)^{-100}.
    \end{equation*}
    Therefore $|\hat{\eta}(\frac{l}{10Q^2})|\lesssim Q^{-100\cdot(S-2)}$. The triangle inequality then gives

    \begin{equation*}
        |\widehat{1-\eta_Q}(l)| \lesssim \log Q\,Q^{-2}Q^{-100\cdot(S-2)} S Q\lesssim \log Q Q^{-1}, \, S\geq 10.
    \end{equation*}
    So in any case, we conclude.
\end{proof}

\end{document}
